\chapter{Conclusion}
In the past 30 years, aerodynamic shape optimization has experienced monumental leaps in its ability to optimize aircraft configurations with incredible detail.
Developments in computational capabilities, numerical optimization methods, and flow solution schemes have enabled optimizations to include an ever-increasing number of design variables.
Yet, despite the widespread capability to conduct analyses at this level, few have expanded beyond traditional wing optimization design variables such as section shape, angle of attack, and twist.
Most importantly, despite the importance of the wing sweep angle in wave drag at typical transonic cruise speeds, its behaviour with relaxed constraints in a pure aerodynamic shape optimization case is unknown.
This study sought to reduce this uncertainty by characterizing sweep and performance changes during ASO of the NASA CRM and A320neo wings with wing sweep as a design variable, and by characterizing the trend in drag benefits as the sweep angle is gradually increased.
It was observed that the performance trends with increasing sweep limits are extremely similar between the NASA CRM and Airbus A320neo wings. Both configurations settled at an optima between 45-50\degree. It was observed that drag benefits can be largely obtained by conducting optimization at the original sweep angle, such that the drag reduction as a result of an increased sweep angle would produce an improvement of only 1.141\% and 2.827\% for the NASA CRM and A320neo, respectively.

The behaviour in drag reduction resulting from increases in the wing sweep angle was studied through ASO on the NASA CRM wing, which was optimized at its original sweep angle of 37.5\degree, along with 40.0\degree and 42.5\degree. This was compared with the drag for the same wing in its original configuration trimmed to the same $C_L S$ constraint, along with the optimized CRM wing at the sweep angle optimum of 48.5\degree. It was observed that initial decreases in drag are sub-linear for small increases in the wing sweep angle, and that further increases in the sweep angle result in continuously diminishing returns. An approximation of the wing weight as a result of the increased wing sweep angle was produced to enable an analysis of net benefits when considering other aircraft design disciplines. It was found that despite the initial appearance of drag benefits with increasing sweep angle, the weight penalty and the subsequent increase in lift-induced drag will negate the drag benefits of increase the sweep angle. This illustrated that increases in the wing sweep angle beyond the nominal value do not yield any net benefits in aircraft performance.

This study shone a light on an approach to aircraft shape optimization that was rarely studied. In doing so, it illustrated the strength of modern ASO methods in conducting optimization with many design variables, along with their ability to exploit the design space in an extremely effective manner. The complexity and subtlety of the interactions between an aircraft's wing sweep angle other aircraft design disciplines cannot be understated. It is recommended that future studies on wing sweep in aircraft design conduct a deeper analysis on the magnitude and cause of drag reduction when the sweep angle approaches the optimum. Furthermore,  Such interactions will undoubtedly increase the complexity of studies on novel technologies with sweep angle considerations, but the result of these efforts will enable more informative conclusions on the path towards improved aircraft efficiency and sustainability in years to come.